% ====================================================================

% ====================================================================
\documentclass[article,paper=a4paper]{jlreq}

% jpreportパッケージの読み込み
\usepackage{jpreport}

% ページレイアウトの設定
\usepackage{geometry}
\usepackage{enumitem}
\geometry{
    top=25mm,
    bottom=25mm,
    left=25mm,
    right=25mm
}
% TikZライブラリ
\usetikzlibrary{
    patterns,
    calc,
    arrows,
    decorations,
    angles,
    shapes.geometric,
    fit,
    knots
}
% TikZ用のスタイル
\tikzset{
  every path/.style={
    black,
    line width=1pt
  },
  every node/.style={
    transform shape,
    knot crossing,
    inner sep=1.5pt
  }
}
% ハイパーリンクの設定
\usepackage[unicode]{hyperref}
\usepackage{bookmark}
\usepackage[nameinlink]{cleveref}

\hypersetup{
    colorlinks=true,
    linkcolor=black,
    urlcolor=blue,
    citecolor=blue,
    pdftitle=,
    pdfauthor=
}

% 数学演算子の定義


% cleveref用の日本語設定
\crefname{thm}{定理}{定理}
\crefname{lem}{補題}{補題}
\crefname{cor}{系}{系}
\crefname{ex}{例}{例}
\crefname{def}{定義}{定義}
\crefname{prop}{命題}{命題}
\crefname{prob}{問題}{問題}

% 色彩の定義
\definecolor{yukari}{RGB}{196, 163, 191}

% 文書情報
\title{代数学 1}
\author{akari0koutya}
\date{\today}

\begin{document}

% --- タイトル ------------------------------------------------------
\maketitle
\tableofcontents
%% 本文の内容をここに記述
% ====================================================================
\section{第1回 : 2026/4/13}
\subsection{ガイダンス}
\section{第2回 : 2026/4/20}
\subsection{復習}
\begin{dfn}
\textbf{アーベル群}(\textbf{可換群}) とは、
\begin{align*}
  \begin{cases}
    \text{集合} \quad G \\
    \text{写像} \quad + \colon G \times G \to G, \quad (a, b) \mapsto a + b 
  \end{cases}
\end{align*}
の組 $(G, +)$ であって、以下の条件を満たすものをいう :
\begin{enumerate}
  \item 結合法則 : $\forall a, b, c \in G, (a + b) + c = a + (b + c)$
  \item 単位元の存在 : $\exists 0 \in G, \forall a \in G, 0 + a = a + 0 = a$
  \item 逆元の存在 : $\forall a \in G, \exists -a \in G, a + (-a) = (-a) + a = 0$
  \item 可換性 : $\forall a, b \in G, a + b = b + a$
\end{enumerate}
\end{dfn}

\begin{rem}
\begin{enumerate}
  \item 組 $(G, +)$ という表記を省略して、単に $G$ と書くこともある。
  \item 単位元 $0$ は一意である。すなわち、もし $0'$ も単位元であれば、$0' = 0$ となる。実際、$0' = 0' + 0 = 0 + 0' = 0$ である。
  \item 逆元 $a^{-1}$ も一意である。実際、 $b$ も $a$ の逆元であれば、
\begin{align*}
  b &= b + 0 \quad (\because 0 \text{ は単位元})\\
    &= b + ((-a) + a) \quad (\because -a \text{ は } a \text{ の逆元})\\
    &= (b + (-a)) + a \quad (\because \text{結合法則})\\
    &= 0 + a \quad (\because b \text{ は } a \text{ の逆元})\\
    &= a
\end{align*}
となり、$b = a^{-1}$ である。
\end{enumerate}
\end{rem}

\begin{dfn}
  \textbf{環} とは、
  \begin{align*}
    \begin{cases}
      \text{集合} \quad R \\
      \text{写像} \quad + \colon R \times R \to R, \quad (a, b) \mapsto a + b \\
      \text{写像} \quad \cdot \colon R \times R \to R, \quad (a, b) \mapsto a \cdot b
    \end{cases}
  \end{align*}
  の組 $(R, +, \cdot)$ であって、以下の条件を満たすものをいう : 
  \begin{enumerate}
    \item $(R, +)$ はアーベル群である。
    \item 結合法則 : $\forall a, b, c \in R, (a \cdot b) \cdot c = a \cdot (b \cdot c)$
    \item 分配法則 : $\forall a, b, c \in R, a \cdot (b + c) = a \cdot b + a \cdot c$ および $(a + b) \cdot c = a \cdot c + b \cdot c$
    \item $\exists 1 \in R, \forall a \in R, 1 \cdot a = a \cdot 1 = a$ 
  \end{enumerate}
  さらに、以下の条件を満たすとき、$R$ は \textbf{可換環} であるという :
\begin{enumerate}[resume]
  \item 可換性 : $\forall a, b \in R, a \cdot b = b \cdot a$
\end{enumerate}
\end{dfn}

\begin{rem}
  特に断らない限り、この資料では、「環」は「可換環」であるとする。
\end{rem}

\begin{rem}
  \begin{enumerate}
    \item 組 $(R, +, \cdot)$ という表記を省略して、単に $R$ と書くこともある。
    \item 環 $R$ の単位元 $1$ は一意である。実際、もし $1'$ も単位元であれば、
    \[1' = 1' \cdot 1 = 1 \cdot 1' = 1\]
    となり、$1' = 1$ である。
  \end{enumerate}
\end{rem}

\begin{dfn}
  \textbf{体} とは、可換環 $(F, +, \cdot)$ であって、以下の条件を満たすものをいう :
  \begin{enumerate}
    \item $F \neq \{0\}$
    \item $\forall a \in F \setminus \{0\}, \exists a^{-1} \in F, a \cdot a^{-1} = a^{-1} \cdot a = 1$
  \end{enumerate}
\end{dfn}

\begin{ex}
  \begin{enumerate}
    \item 次は体である :
      \begin{align*}
        \begin{cases}
          \Q : \text{有理数全体} \\
          \R : \text{実数全体} \\
          \C : \text{複素数全体}
        \end{cases}
      \end{align*}
    \item 次は環であるが体ではない :
      \begin{align*}
        \begin{cases}
          \Z : \text{整数全体} \\
          \C[x] : \C \text{係数の多項式全体}\\
          K[x] : K \text{係数の多項式全体} \quad (K \text{ は体})
        \end{cases}
      \end{align*}
    \item $\Z[\sqrt{-1}] = \{a + b\sqrt{-1} \mid a, b \in \Z\}$ は環である。実際、$a+ b\sqrt{-1}, c + d\sqrt{-1} \in \Z[\sqrt{-1}]$ に対し、
      \begin{align*}
        (a + b\sqrt{-1}) + (c +  d\sqrt{-1}) = (a + c) + (b + d)\sqrt{-1} &\in \Z[\sqrt{-1}] \\
        (a + b\sqrt{-1})(c + d\sqrt{-1}) = (ac - bd) + (ad + bc)\sqrt{-1} &\in \Z[\sqrt{-1}]
      \end{align*}
      なので、$\Z[\sqrt{-1}]$ は環である。しかし、$2 \in \Z[\sqrt{-1}]$ に対して、$2^{-1} = \frac{1}{2} \not\in \Z[\sqrt{-1}]$ なので、$\Z[\sqrt{-1}]$ は体ではない。同様にして、$\Q(\sqrt{-1})$ も環であるが、これは体である。実際、$a+b\sqrt{-1} \in \Q(\sqrt{-1})$ に対して、
      \[(a+b\sqrt{-1})^{-1} = \frac{a-b\sqrt{-1}}{a^2+b^2} = \frac{a}{a^2+b^2} - \frac{b}{a^2+b^2}\sqrt{-1} \in \Q(\sqrt{-1})\] 
      なので、$\Q(\sqrt{-1})$ は体である。
    \item $n \geq 2$ としたとき、$\Z/n\Z$ は環である。実際、$a + n\Z, b + n\Z \in \Z/n\Z$ に対し、
      \begin{align*}
        (a + n\Z) + (b + n\Z) = (a + b) + n\Z &\in \Z/n\Z \\
        (a + n\Z)(b + n\Z) = ab + n\Z &\in \Z/n\Z
      \end{align*}
      なので、$\Z/n\Z$ は環である。\\
      \quad $n$ が合成数のとき、$\Z/n\Z$ は体ではない。実際、$n = ab$ としたとき、
      \[(a + n\Z)(b + n\Z) = 0 + n\Z \in \Z/n\Z\] 
      となるので、$\Z/n\Z$ は体ではない。\\
      \quad 一方、$\Z/n\Z$ が体であるのは、$n$ が素数のときである。実際、$\Z/n\Z$ が体であれば、$\exists k \in \Z$ に対して、
      \[(k + n\Z)(a + n\Z) = ka + n\Z =  1 + n\Z.\]
      $ka + n\Z$ という集合には、$a$ の倍数しか属さないので、$ka + n\Z = 1 + n\Z$ となるためには、$a$ と $n$ が互いに素でなければならない。したがって、$\Z/n\Z$ が体であるためには、$n$ が素数でなければならない。\\
      \quad 例えば、$\Z/5\Z = \{\overline{0} + 5\Z, \overline{1} + 5\Z, \overline{2} + 5\Z, \overline{3} + 5\Z, \overline{4} + 5\Z\}$ は体である。実際、
      \begin{align*}
        \begin{cases}
          (\overline{1} + 5\Z)(\overline{1} + 5\Z) = 1 + 5\Z \quad &\therefore (\overline{1} + 5\Z)^{-1} = \overline{1} + 5\Z \\
          (\overline{2} + 5\Z)(\overline{3} + 5\Z) = 6 + 5\Z = 1 + 5\Z \quad &\therefore (\overline{2} + 5\Z)^{-1} = \overline{3} + 5\Z \\
          (\overline{3} + 5\Z)(\overline{2} + 5\Z) = 6 + 5\Z = 1 + 5\Z \quad &\therefore (\overline{3} + 5\Z)^{-1} = \overline{2} + 5\Z \\
          (\overline{4} + 5\Z)(\overline{4} + 5\Z) = 16 + 5\Z = 1 + 5\Z \quad &\therefore (\overline{4} + 5\Z)^{-1} = \overline{4} + 5\Z \\
        \end{cases}
      \end{align*}
      なので、$\Z/5\Z$ は体である。\\
      \quad $n$ が一般の素数であるとし、$m + n\Z \in \Z/n\Z (m \in n\Z)$ とする。このとき、 $m$ と $n$ は互いに素であるので、$\exists k, l \in \Z$ に対して、$km + ln = 1$ となる。したがって、
      \[(k + n\Z)(m + n\Z) = km + n\Z = 1 + n\Z\]
      となるので、$(m + n\Z)^{-1} = k + n\Z$ である。したがって、$\Z/n\Z$ は体である。
  \end{enumerate}
\end{ex}

\subsection{環上の加群}

\begin{dfn}
  環 $R$ に対し、R\textbf{-加群} とは、
  \begin{align*}
    \begin{cases}
      \text{集合} \quad M \\
      \text{写像} \quad + \colon M \times M \to M, \quad (x, y) \mapsto x + y \\
      \text{写像} \quad R \times M \to M, \quad (r, x) \mapsto rx
    \end{cases}
  \end{align*}
  という組 $(M, +, R \times M \to M)$ であって、以下の条件を満たすものをいう :
\begin{enumerate}
  \item $(M, +)$ はアーベル群
  \item $\forall r, s \in R$, $\forall x \in M$, $(rs)x = r(sx)$
  \item $\forall r, s \in R$, $\forall x \in M$, $(r + s)x = rx + sx$
  \item $\forall r \in R$, $\forall x, y \in M$, $r(x + y) = rx + ry$
  \item $\forall x \in M$, $1x = x$
\end{enumerate}
\end{dfn}

\begin{rem}
  $R$ が体 $K$ であるとき、$K$-加群とは\textbf{ベクトル空間}のことである。
\end{rem}

\clearpage

\begin{ex}
  $n \geq 1$ のとき、
  \[\R^n = \left\{ \begin{pmatrix} a_1 \\ a_2 \\ \vdots \\ a_n \end{pmatrix} \relmid | a_i \in \R \right\}\]
  は $\R$-加群である。実際、$x = \begin{pmatrix} a_1 \\ a_2 \\ \vdots \\ a_n \end{pmatrix}, y = \begin{pmatrix} b_1 \\ b_2 \\ \vdots \\ b_n \end{pmatrix} \in \R^n$ と $r \in \R$ に対し、
\begin{align*}
  x + y &= \begin{pmatrix} a_1 + b_1 \\ a_2 + b_2 \\ \vdots \\ a_n + b_n \end{pmatrix} \in \R^n \\
  rx &= \begin{pmatrix} ra_1 \\ ra_2 \\ \vdots \\ ra_n \end{pmatrix} \in \R^n
\end{align*}
なので、$\R^n$ は $\R$-加群である。\\
\indent $\R^0 = \{0\}$ とする。
\end{ex}

\section{第3回 : 2026/4/27}

\begin{exe}
  \quad $K$ : 体, $a \in K\backslash\{0\}$.\\
  このとき, $ab = 1$ ならば $b = a^{-1}$ であることを示せ。
\end{exe}
\begin{proof}[解答]
  $ab = 1$ と仮定するこのとき、
  \begin{align*}
    a^{-1} &= a^{-1}\cdot 1\\
    &= a^{-1} \cdot (ab) \\
    &= (a^{-1} \cdot a) \cdot b \\
    &= 1 \cdot b \\
    &= b
  \end{align*}

  より, $b = a^{-1}$ である。
\end{proof}

\begin{exe}
  \quad $M$ : $R$-加群, $x \in M$.\\
  このとき, $0x = 0$ であることを示せ。
\end{exe}

\begin{proof}[解答]
  分配法則より, 
  $0x = (0 + 0)x = 0x + 0x$ であるので、両辺から $0x$ を引くと $0 = 0x$ となる。
\end{proof}

\begin{ex}
  アーベル群は自然に $\Z$-\hspace{0pt}加群と見做せる。\\
  $M$ : アーベル群, $x \in M$\\
  このとき, 
  \[2\cdot x = (1+1)\cdot x = 1\cdot x + 1\cdot x = x+x.\]
  このように、$n \in \Z$ に対して 
  \[n\cdot x = \begin{cases}
  \underbrace{x + x + \cdots + x}_{n \text{ 個}} & (n > 0) \\
  0 & (n = 0) \\
  \underbrace{(-x) + (-x) + \cdots + (-x))}_{n \text{ 個}} & (n < 0)
  \end{cases}\]
  とすると、$M$ は $\Z$-\hspace{0pt}加群と見做せることがわかる。
\end{ex}

\begin{rem}
  \quad $R$ : 環, $M$ : $R$-加群.\\
  \begin{enumerate}
    \item $\forall x\in M$, $0x = 0$.
    \item $\forall x\in M$, $\underbrace{(-1)x}_{-1 \in R, x \in M} = \underbrace{-x}_{-x\in M : M\text{における}x\text{の下方に関しての逆元}}$.
  \end{enumerate}
  一つ目に関しては、すでに示した通りである。二つ目に関しては、
  \begin{align*}
    (-1)x + x &= (-1)x + 1x \\
    &= (-1 + 1)x \\
    &= 0x \\
    &= 0
  \end{align*}

  $\Z/2\Z = \{2\Z, \overline{1}+2\Z\}$ はアーベル群でり、上の例より、$\Z$-\hspace{0pt}加群と見做せる。これは、$\Z^n$ と同型でない。実際、$\Z^n$ は$\{0\}$ あるいは無限集合であるのに対し、$\Z/2\Z$ は$2$元集合である。
\end{rem}

\clearpage

\begin{dfn}
  \quad $R$ : 環\\
  $R$の\textbf{イデアル}とは、$I \subseteq R$ であって、以下の条件を満たすものをいう :
\begin{enumerate}
  \item $\forall a, b \in I, a + b \in I$
  \item $\forall a \in I, \forall r \in R, ra \in I$
\end{enumerate}
\end{dfn}

\begin{ex}
  $R$ のイデアル $I$ は $R$-加群である。\\
  実際、$a, b \in I$ と $r \in R$ に対し、
  \[a\in I \Rightarrow 1a=(-1)\cdot a\in I.\]
  また、$R/I$ も $R$-加群である。実際、$a + I, b + I \in R/I$ と $r \in R$ に対し、
  \begin{align*}
    (a + I) + (b + I) &= (a + b) + I \in R/I \\
    r(a + I) &= ra + I \in R/I
  \end{align*}
\end{ex}

\subsection{部分加群と剰余加群}

\begin{dfn}
  \quad $R$ : 環, $M$ : $R$-加群\\
  $M$の\textbf{部分加群}とは、$N \subseteq M$ であって、以下の条件を満たすものをいう :
\begin{enumerate}
  \item $\forall x, y \in N, x + y \in N$
  \item $\forall x \in N, \forall r \in R, rx \in N$
\end{enumerate} 
\end{dfn}

\begin{rem}
  $R$ 自身を $R$-加群と見做すと、$R$ のイデアルは $R$ の部分加群である。
\end{rem}

\begin{dfn}
  \quad $R$ : 環, $M$ : $R$-加群, $N$ : $M$ の部分加群\\
  このとき、アーベル群としての剰余群 $M/N = \{x + N \mid x \in M\}$ に対して、$R$-加群構造が次のように定まる :
\begin{align*}
  (x + N) + (y + N) &= (x + y) + N \\
  r(x + N) &= rx + N
\end{align*}
このように定めると、$M/N$ は $R$-加群となる。これは well-definedに定まっている。これを\textbf{剰余加群}という。
\end{dfn}

\begin{rem}
  $M/N$ とはなんだろうか？\\
  $x \in M$ に対して、
  \[x+N = \{x+n \mid n\in N\} \]
  と定め、
  \[M/N := \{x+N \mid x\in M\}\]
  する。ここで、$x, y \in M$ に対して、
  \[x \sim y \iff x - y \in N\]
  と定めと、この $\sim$ は $M$ 上の同値関係である。
  この同値関係により、$M$ を\textbf{類別}することができる。つまり、$M = \bigcup_{\lambda \in \Lambda} M_{\lambda} $ であって、
  \begin{align*}
    \begin{cases}
      M_{\lambda} \cap M_{\mu} = \emptyset & (\lambda \neq \mu) \\
      M_{\lambda} = \{x \in M \mid x \sim x_{\lambda}\} & (\lambda \in \Lambda)
    \end{cases}
  \end{align*}
  各 $M_{\lambda}$ は、$x_{\lambda}$ を用いて $x_{\lambda}+N$ と表せることが証明できる。
\end{rem}

\begin{rem}
  \quad $M/N$ の演算の well-definedness について、つまり、
  \[(x+N)+(y+N)=(x+y)+N\]
  が矛盾なく"定義"できているか、を吟味する。\\
  $x, y, x', y' \in M$ として、
  \begin{align*}
    \begin{cases}
      x+N = x^{\prime}+N \\
      y+N = y^{\prime}+N
    \end{cases}
  \end{align*}
  とすると、$x - x^{\prime} \in N$ かつ $y - y^{\prime} \in N$ である。したがって、
  \[(x+y) - (x^{\prime} + y^{\prime}) = (x - x^{\prime}) + (y - y^{\prime}) \in N\]
  なので、$(x+y)+N = (x^{\prime} + y^{\prime}) + N$  である。\\
  \quad 同様にして、$r(x+N) = rx + N$ も well-defined に定まっていることがわかる。
  このことから、$M/N$ の加法の定義は一見、大表現の取り方に依存しているが、剰余類の選び方にしか依存していないことがわかる。\\
  \quad 尚、 $M/N$ の $0$ は $0+N$ である。\\
  実際、
  \[(0+N)+(m+N) = (0+m)+N=m+N\]
  なので、$0+N$ は $M/N$ の単位元である。
\end{rem}

\section{第4回 : 2026/5/01}

\begin{rem}
  $M$ : $R$-\hspace{0pt}加群, $N$ : $M$ の部分$R$-\hspace{0pt}加群\\
  このとき、$M/N = \{x+N | x\in M\}$ に対して、$a\in R, x+N\in M/N$ に対して、スカラー倍を
  \[a(x+N) = ax + N\]
  と定義したい。この定義は、$M/N$ の加法の定義と同様に、well-defined に定まっていることを確認する必要がある。\\
  $a\in R, x, y \in M$ として、
  \begin{align*}
    \begin{cases}
      x+N = x^{\prime}+N \\
      y+N = y^{\prime}+N
    \end{cases}
  \end{align*}
  とすると、$x - x^{\prime} \in N$ かつ $y - y^{\prime} \in N$ である。したがって、 
  \begin{align*}
    a(x+y) - a(x^{\prime} + y^{\prime}) &= ax + ay - ax^{\prime} - ay^{\prime} \\
    &= a(x - x^{\prime}) + a(y - y^{\prime}) \in N
  \end{align*}
  なので、$a(x+y) + N = a(x^{\prime} + y^{\prime}) + N$ である。\\
  \quad 同様にして、$a(x+y) + N = a(x^{\prime} + y^{\prime}) + N$ も well-defined に定まっていることがわかる。

  また、$M/N$ が $R$-\hspace{0pt}加群であることは容易にわかる。例えば、結合法則は、
  \begin{align*}
    a(b(x+N)) &= a(bx + N) \\
    &= (a(bx)) + N \\
    &= ((ab)x) + N \\
    &= (ab)(x+N)
  \end{align*}
  である。そのほかは、$M$ が $R$-\hspace{0pt}加群であることから、同様にして示せる。
\end{rem}

\subsection{準同型写像と準同型定理}

\begin{dfn}
  \quad $R$ : 環, $M, N$ : $R$-加群\\
  このとき、写像 $f \colon M \to N$ が$R$\textbf{-\hspace{0pt}線形写像}または$R$\textbf{-\hspace{0pt}加群の準同型写像}であるとは、$\forall x, y \in M, \forall r \in R$ に対して、
  \begin{align*}
    f(x+y) &= f(x) + f(y) \\
    f(rx) &= rf(x)
  \end{align*}
  を満たすことをいう。
\end{dfn}

\begin{dfn}
  \quad $R$ : 環, $M, N$ : $R$-加群\\
  このとき、$M$と$N$が\textbf{同型}であるとは、$R$-\hspace{0pt}線形写像 $f \colon M \to N$ と $g \colon N \to M$ であって、
  \[g \circ f = \mathrm{id}_M, \quad f \circ g = \mathrm{id}_N\]
  を満たすものが存在することをいう。また、このとき、$f$ は同型（写像）であるという。
\end{dfn}

\begin{prop}
  \quad $R$ : 環, $M, N$ : $R$-加群\\
  このとき、次のの条件は同値である :
\begin{enumerate}[label=(\arabic*)\,]
  \item $M$ と $N$ は同型である。
  \item $M$ と $N$ の間に全単射な $R$-\hspace{0pt}線形写像が存在する。
\end{enumerate}
\end{prop}

\begin{proof}[(1) $\Rightarrow$ (2)]
  $M$ と $N$ は同型であると仮定する。このとき、$R$-\hspace{0pt}線形写像 $f \colon M \to N$ と $g \colon N \to M$ であって、
  \[g \circ f = \mathrm{id}_M, \quad f \circ g = \mathrm{id}_N\]
  を満たすものが存在する。特に、$f$ は全単射である。
\end{proof}

\begin{proof}[(2) $\Rightarrow$ (1)]
  $M$ と $N$ の間に全単射な $R$-\hspace{0pt}線形写像 $f \colon M \to N$ が存在すると仮定する。このとき、$f$ は全単射であるので、逆写像 $f^{-1} \colon N \to M$ が存在する。さらに、$f^{-1}$ は $R$-\hspace{0pt}線形写像であることがわかる。実際、$y_1, y_2 \in N$ と $r \in R$ に対し、
  \begin{align*}
    f^{-1}(y_1 + y_2) &= f^{-1}(y_1) + f^{-1}(y_2) \\
    f^{-1}(ry_1) &= f^{-1}(r f(f^{-1}(y_1))) \qquad (\because f \text{ は逆写像をもつ})\\
    &= f^{-1}(f(r f^{-1}(y_1))) \qquad (\because f \text{ は } R\text{-線形})\\
    &= r f^{-1}(y_1)
  \end{align*}
  である。したがって、$f$ と $f^{-1}$ は同型な写像であるので、$M$ と $N$ は同型である。
\end{proof}

\begin{dfn}
  \quad $R$ : 環, $M, N$ : $R$-加群, $f \colon M \to N$ : $R$-\hspace{0pt}線形写像
  \begin{enumerate}[label=(\arabic*)\,]
    \item $f$ の\textbf{核}とは、$\Ker(f) := \{x \in M \mid f(x) = 0\}$ をいう。
    \item $f$ の\textbf{像}とは、$\im(f) := \{f(x) \mid x \in M\}$ をいう。
  \end{enumerate}
\end{dfn}

\begin{prop}
  \quad $R$ : 環, $M, N$ : $R$-加群, $f \colon M \to N$ : $R$-\hspace{0pt}線形写像\\
  このとき、以下が成り立つ :
\begin{enumerate}[label=(\arabic*)\,]
  \item $\Ker(f)$ は $M$ の部分$R$-\hspace{0pt}加群である。
  \item $\im(f)$ は $N$ の部分$R$-\hspace{0pt}加群である。
\end{enumerate}
\end{prop}

\begin{proof}[(1)]
  $x, y \in \Ker(f)$ と $r \in R$ に対し、
  \begin{align*}
    f(x+y) &= f(x) + f(y) = 0 + 0 = 0 \\
    f(rx) &= rf(x) = r\cdot 0 = 0
  \end{align*}
  なので、$x+y \in \Ker(f)$ かつ $rx \in \Ker(f)$ である。したがって、$\Ker(f)$ は $M$ の部分$R$-\hspace{0pt}加群である。
\end{proof}

\begin{proof}[(2)]
  $y_1, y_2 \in \im(f)$ と $r \in R$ に対し、それぞれ $x_1, x_2 \in M$ が存在して $f(x_1) = y_1, f(x_2) = y_2$ である。このとき、
  \begin{align*}
    y_1 + y_2 &= f(x_1) + f(x_2) = f(x_1 + x_2) \\
    ry_1 &= r f(x_1) = f(rx_1)
  \end{align*}
  なので、$y_1 + y_2 \in \im(f)$ かつ $ry_1 \in \im(f)$ である。したがって、$\im(f)$ は $N$ の部分$R$-\hspace{0pt}加群である。
\end{proof}

\begin{rem}
  重要な性質として、以下のものがある :
  \[f(0) = f(0+0) = f(0) + f(0)\]
  したがって、$f(0) = 0$ である。
\end{rem}

\begin{thm}
  \quad $R$ : 環, $f \colon M \to N$ : $R$-\hspace{0pt}線形写像\\
  \quad $\pi \colon M \to M/\Ker(f), x \mapsto x + \Ker(f)$ \\
  このとき、$R$-\hspace{0pt}線形写像 $g \colon M/\Ker(f) \to N$ がただ一つ存在し、
  \[g(x + \Ker(f)) = f(x)\]
  が成り立つ。また、この $g$ は同型
  \[M/\Ker(f) \overset{\cong}{\to} \im(f)\]
  を定める。\\
  $g$ の条件は次のように表すこともできる :\\
  \begin{center}
    \begin{tikzcd}
  A \ar[r, "f"] \arrow[d, "g"'] & B \ar[d] \\
  C \ar[r] & D
\end{tikzcd}
  \end{center}
\end{thm}

\begin{proof}
  $g$ が存在することを示す。$x + \Ker(f) \in M/\Ker(f)$ に対して、$g(x + \Ker(f)) := f(x)$ と定める。これが well-defined に定まっていることを確認する。$x, y \in M$ として、
  \begin{align*}
    \begin{cases}
      x + \Ker(f) = y + \Ker(f) \\
      f(x) = g(x + \Ker(f)) \\
      f(y) = g(y + \Ker(f))
    \end{cases}
  \end{align*}
  とすると、$x - y \in \Ker(f)$ であるので、$f(x - y) = 0$ である。したがって、$f(x) = f(y)$ である。したがって、$g(x + \Ker(f)) = g(y + \Ker(f))$ である。\\

  \vspace{1\baselineskip}

  \quad $g$ は $R$-\hspace{0pt}線形写像であることを確認する。$x, y \in M$ と $r \in R$ に対し、
  \begin{align*}
    g((x+y) + \Ker(f)) &= g((x+\Ker(f)) + (y+\Ker(f)))\\
    &= g(x+\Ker(f)) + g(y+\Ker(f))\\
    &= f(x) + f(y)\\
    &= f(x+y)
  \end{align*}
  および
  \begin{align*}
    g(r(x+\Ker(f))) &= g(rx + \Ker(f))\\
    &= f(rx)\\
    &= rf(x)
  \end{align*}
  より、$g$ は $R$-\hspace{0pt}線形写像である。\\

  \vspace{1\baselineskip}

  $g$ は一意的であることを確認する。$g' \colon M/\Ker(f) \to N$ も $R$-\hspace{0pt}線形写像であって、$g'(x + \Ker(f)) = f(x)$ を満たすとする。このとき、$\forall x + \Ker(f) \in M/\Ker(f), g(x + \Ker(f)) = g'(x + \Ker(f))$ である。したがって、$g$ は一意的である。\\

  \vspace{1\baselineskip}

  $g$ は同型であることを確認するために、全単射であることを確認する。\\
  \vspace{1\baselineskip}
  \quad 全射であることは、$\forall y\in\im(f),\exists x\in M/\Ker(f)$ に対して、$g(x)=y$ を示せば十分である。$\forall y\in\im(f),\exists x_0\in M$ に対して、$f(x_0)=y$ である。このとき、
  \[g(x_0 + \Ker(f)) = f(x_0) = y\]
  なので、$g$ は全射である。\\
  \quad 単射であることは、$\forall x\in M/\Ker(f), g(x)=0 \Rightarrow x=0$ を示せば十分である。$x\in M/\Ker(f)$ として、$g(x)=0$ と仮定する。このとき、$x = x_0 + \Ker(f)$ と表せる $x_0 \in M$ が存在する。このとき、
  \[0 = g(x) = g(x_0 + \Ker(f)) = f(x_0)\]
  なので、$x_0 \in \Ker(f)$ である。したがって、$x = x_0 + \Ker(f) = 0 + \Ker(f) = 0$ である。\\
  $g$ は同型であることがわかったので、$g^{-1} \colon \im(f) \to M/\Ker(f)$ が存在する。さらに、$g^{-1}$ は $R$-\hspace{0pt}線形写像であることがわかる。実際、$y_1, y_2 \in \im(f)$ と $r \in R$ に対し、
  \begin{align*}
    g^{-1}(y_1 + y_2) &= g^{-1}(y1) + g^{-1}(y_2) \\
    g^{-1}(ry_1) &= g^{-1}(r g(g^{-1}(y_1))) \qquad (\because g \text{ は逆写像をもつ})\\
    &= g^{-1}(g(r g^{-1}(y_1))) \qquad (\because g \text{ は } R\text{-線形})\\
    &= r g^{-1}(y_1)
  \end{align*}
  である。したがって、$g$ と $g^{-1}$ は同型な写像であるので、$M/\Ker(f)$ と $\im(f)$ は同型である。
\end{proof}

\begin{ex}
  \quad $f(x)$ : $\R$ 上の $C^{\infty}$ 関数, $F(x)$ : $f(x)$ の原始関数, \\
  \quad $C^{\infty}(\R)$ : $\R$ 上の $C^{\infty}$ 関数全体の集合, $D \colon C^{\infty}(\R) \to C^{\infty}(\R), f \mapsto f'$ : $C^{\infty}(\R)$ 上の微分作用素
  \vspace{\zh}
  このとき、$D$ は $C^{\infty}(\R)$ 上の $\R$-\hspace{0pt}線形写像である。さらに、
  \begin{align*}
    \Ker(D) &= \{F(x) \in C^{\infty}(\R) \mid F'(x) = 0\}\\
    &= \{\R \text{上の定数関数}\}\\
    &= \R
  \end{align*}
  また、
  \begin{align*}
    \im(D) &= \{f'(x) \in C^{\infty}(\R) \mid f(x) \in C^{\infty}(\R)\}\\
    &= C^{\infty}(\R)
  \end{align*}
  よって、準同型定理より、\[C^{\infty}(\R)/\R \overset{\cong}{\to} C^{\infty}(\R)\] である。
  $D$ の逆写像は不定積分 $\int f(x) dx$ である。実際、$f(x) \in C^{\infty}(\R)$ に対して、$f(x) + \R \in C^{\infty}(\R)/\R$ と定めると、$C^{\infty}(\R)/\R$ の元である $f(x) + \R$ は $C^{\infty}(\R)$ の元である $\int f(x) dx$ と同一視できる。\\
  たとえば、$f(x) = e^x$ とすると、$f(x) + \R$ は $C^{\infty}(\R)/\R$ の元である。これは不定積分 $\int e^x dx$ と同一視できる。
\end{ex}
\end{document}